\label{chapter_lit_review}
In this chapter, we define the main approaches to solve the problem with  examples for each approach.\\

The problem of clearing traffic for an emergency vehicle was addressed in the literature using three approaches:
\begin{itemize}
    \item \textbf{Vehicle rerouting}, which involves planning alternative routes to clear the congestion or avoid it.
    \item \textbf{Dynamic signal control}, which depends on a control center that directs the traffic flow in a way that enables the emergency vehicle to reach its destination with minimum delay through controlling traffic signals in its area. 
    \item \textbf{Independent vehicle planning}, where vehicles independently search and plan for possible routes for the traffic clearance mission.
\end{itemize}
\section{Vehicle Rerouting}
\label{vehicle_rerouting}
\cite{paper1} suggests informing emergency vehicles ahead of time which roads are the most congested so that they could avoid taking them to their destinations. In \cite{paper20}, vehicles within a certain distance from the incoming emergency vehicle were informed about it so that their drivers could take the action they see appropriate.. On another hand, in \cite{paper30}, each vehicle was not only warned of the approaching emergency vehicle but also instructed, through a mobile application, with the exact suitable action its driver ought to take in order to make way for the emergency vehicle. This warning message contained information about the emergency vehicle's current position, velocity, heading, destination, and route information. 
In \cite{paper4}, a very similar solution was presented. A simulation of one ambulance and two other vehicles was tested and the decision algorithm succeeded in instructing the vehicles with the truly appropriate maneuvers.
\\ \\
These approaches have the merit of solving the emergency traffic clearance problem, however, a major demerit that should be noted is the possibility of not finding a solution or that the found solution might be longer.

\section{Dynamic Signal Control}
\label{Dynamic_signal_control}
Various publications attempted the use of traffic information to generate a sequence of green light signals prioritizing the emergency vehicle. \cite{paper5} proposes establishing communication between traffic lights and emergency vehicles such that they are prioritized at intersections and always given green lights. \cite{paper6} proposed using light emitters that notice problems such as blocked line-of-sight and excessive noise. Similarly, \cite{paper7} used technologies such as IR, and GPS to estimate traffic density and emergency vehicles presence. Others like \cite{paper8} combined measurements of the distance between the emergency vehicle and an intersection using visual sensing methods, vehicle counting and time-sensitive alert transmission within the sensor network to plan for traffic signals sequence. They concluded that using Euclidean distance outperforms other techniques such as Manhattan distance and Canberra distance.
\\ \\ 
These approaches show promising results. However, they have many drawbacks such as not considering non-emergency vehicles’ individual goals as they might be stuck in a red light due to poor planning for the emergency vehicle. Moreover, they are useless if the congestion is due to an accident for instance, and suffer from varying planning strategies in different coverage areas with each having its own intelligent control center.
\section{Independent Vehicle Planning}
\label{independent_learner}
In \cite{paper9}, David L. et al discusses cooperation between autonomous vehicles to perform lane merging with no communication needed in their paper “Tactical Cooperative Planning for Autonomous Highway Driving using Monte-Carlo Tree Search”. The paper starts with presenting different rule-based algorithms such as “finite state machine” that were previously adopted to solve the problem. Then it notes that these algorithms don’t suit the uncertainties in the input.  Hence, they proposed using “Monte-Carlo Tree” MCT search to parse a tree of possible actions taken by each agent to minimize a joint cost while minimizing its own cost. Each agent, named “ego” in the literature, observes an environment of other interacting agents marked in blue, along with vehicles that are influenced by its decision marked in white.
\\ \\
The authors utilized the powerful capabilities of MCT while adding tactical and cooperative motion planning,TCMP, to come up with TCMP-MCT algorithm. In TCMP-MCT, A simulation period between two consecutive actions yields a combined cost function to guide the enhancement of both, the agent’s plan, and the collective plan of the whole system. The main problem with this algorithm is that actions are taken sequentially by the agents in order for them to be organized in a tree form. Such a model should result in a behavior where agents learn to follow the actions of their neighbors rather than executing a joint simultaneous action.
The proposed methodology proved successful in experiments concerned with planning highway driving scenarios that have a partially observable MDP, POMDP, nature. However, it lacks the powerful merit of V2V communication. Vehicles act independently with no cooperation and hency in real actio, their plans might contradict.

\section{Reinforcement Learning as an Alternative}
\label{sec:RL_as_alternative}
In light of the previously discussed literature review and the research gaps identified, we found that RL offers more promising approaches. In recent years, Rl algorithms were capable of competing with the previously dominant supervised learning techniques as they need little to no intervention from humans. In \cite{paper10} for instance, a deep RL approach plays the game of AlphaGo based only on trial and error as an RL learning methodology. Also, the work of \cite{paper11} develops a variant of Q-Learning, deep Q-Learning (DQRL), to play Atari on the Arcade Learning Environment. The algorithm not only outperforms all other existing algorithms in six games played, but it actually outperforms human players in 3 of them. This trial and error nature of RL algorithms makes it less problem defined, and hence easily transferred to other tasks. In addition, RL gives space for multi-agent learning,  MARL, where connected agents can perform even better. \cite{paper12} suggests that independently training single agents might show efficiency in certain tasks, but it’s ultimately optimal to make the agents coordinate together. 
\\ \\ 
In a MARL setting, multiple vehicles can use the availability of V2V communication to cooperate on clearing the road for the vehicle. However,  with cooperation comes many questions; such as centralization and decentralization for instance. Who should give the orders? A central unit, a leading agent, or it’s a form of anarchy where agents independently make decisions affecting others. An answer to this question was found in “Reinforcement Learning State-of-the-Art” by M. V. Otterlo \cite{rl_state_of_the_art}, where the problem is defined as an MDP game. The book proposes using value iteration based algorithms that could help arrive at nash equilibrium between players (agents) such as joint action learner (JAL). In JAL, a joint action is taken based on observing both the environment and a probabilistic estimation of the possible actions made by the agents learning from previous experience of such state-joint action encounters. A noteworthy work on MARL, \cite{paper13}, uses a DQRL approach with a value decomposition network. Thus, Q function estimation would depend on the addition of all agent’s Q functions to relieve the problem of conflicting behaviors or lazy agents. 


